\documentclass[a4paper,10pt]{article}
\usepackage{amsmath,amssymb}
\usepackage{fontspec}
\setmainfont{Liberation Serif}
\usepackage{graphicx}
\usepackage{eso-pic}
\usepackage{tikz}
\usepackage[most]{tcolorbox}
\usepackage{geometry}
\usepackage{fancyhdr}
\usepackage{lastpage}
\usepackage{booktabs}
\usepackage{listings}
\usepackage{enumitem}
\usepackage{microtype}

\geometry{
    a4paper,
    top=1.5cm,
    bottom=1.8cm,
    left=1.2cm,
    right=1.2cm
}

% Persistent Unremovable Canvas Watermark on Every Page
\newsavebox{\watermarkbox}
\sbox{\watermarkbox}{%
    \begin{tabular}{c}
        \includegraphics[width=8cm]{/home/johan/GATE/GATE_CS_PYQ_DB/assets/studiva-butterfly-icon.jpg} \\[-0.5em]
        {\fontsize{18}{22}\selectfont\bfseries\scshape www.studiva.co.in}
    \end{tabular}%
}
\AddToShipoutPictureBG{%
    \begin{tikzpicture}[remember picture, overlay]
        \node[opacity=0.12] at (current page.center) {\usebox{\watermarkbox}};
    \end{tikzpicture}%
}

% Header and Footer
\pagestyle{fancy}
\fancyhf{}
\fancyhead[L]{\small\scshape GATE 2018--2025 CS/IT $\cdot$ \textbf{Databases}}
\fancyhead[R]{\small\scshape Practice Question Bank with Solutions}
\fancyfoot[L]{\scriptsize\textsc{Copyright \copyright\ 2026} Studiva\textsuperscript{TM}}
\fancyfoot[C]{\small Page \thepage\ of \pageref{LastPage}}
\fancyfoot[R]{\scriptsize\texttt{www.studiva.co.in}}
\renewcommand{\headrulewidth}{0.4pt}
\renewcommand{\footrulewidth}{0.4pt}

% Code Listing Style (Monochrome)
\lstset{
    basicstyle=\footnotesize\ttfamily,
    numbers=none,
    frame=none,
    tabsize=2,
    breaklines=true,
    showstringspaces=false
}

\begin{document}

% ==================== PAGE 1: IMPRINT & TITLE PAGE ====================
\thispagestyle{empty}
\vspace*{0.5cm}
\begin{center}
    \includegraphics[width=2.5cm]{/home/johan/GATE/GATE_CS_PYQ_DB/assets/studiva-butterfly-icon.jpg}\\[0.6em]
    {\LARGE \textbf{GATE 2018--2025 COMPUTER SCIENCE \& IT}}\\[0.3em]
    {\Large \textbf{PREVIOUS YEAR QUESTION BANK WITH SOLUTIONS}}\\[0.4em]
    {\large \textsc{\textbf{Subject Wise: Databases}}}\\[0.2em]
    {\normalsize \textsc{Combined 2018--2025 Examination Papers $\cdot$ Comprehensive Edition ($\mathbf{31}$ Questions)}}\\[1.2em]
    \rule{0.6\textwidth}{0.6pt}
\end{center}

\vspace{1em}
\noindent\begin{tabular}{ll}
\textbf{Author / Compiler:} & Rupam Ghosh \\
\textbf{Publisher \& Platform:} & Studiva\textsuperscript{TM} (\texttt{www.studiva.co.in}) \\
\textbf{Official Inquiries:} & \texttt{contact@studiva.co.in} \\
\textbf{Edition:} & First Standard Edition (2026) \\
\textbf{Coverage:} & Complete GATE CS/IT Papers (2018, 2022, 2023, 2024 CS1, 2024 CS2, 2025 CS1, 2025 CS2) \\
\textbf{Syllabus Reference:} & Official GATE 2026 Syllabus (Organizing Institute: IIT Guwahati) \\
\end{tabular}

\vspace{1.5em}
\begin{tcolorbox}[enhanced, standard jigsaw, opacityback=0, colframe=black!40, arc=1mm, boxrule=0.5pt, left=10pt, right=10pt, top=8pt, bottom=8pt]
\textbf{\large COPYRIGHT \& INTELLECTUAL PROPERTY NOTICE}\\[0.6em]
Copyright \copyright\ 2026 Rupam Ghosh \& Studiva\textsuperscript{TM}. All Rights Reserved.\\[0.5em]
This educational compilation, problem organization, taxonomy classification, LaTeX typesetting, hints, and explanatory solutions are the proprietary intellectual property of \textbf{Rupam Ghosh} and \textbf{Studiva\textsuperscript{TM}}.\\[0.6em]
\textbf{Terms of Use \& Attribution Rules:}
\begin{itemize}[leftmargin=*, nosep]
    \item \textbf{Personal Educational Use:} Permitted freely for individual students and educators preparing for competitive examinations.
    \item \textbf{Mandatory Attribution:} Any reproduction, republication, hosting, or sharing (digital or physical) must prominently credit \textbf{Rupam Ghosh \& Studiva\textsuperscript{TM}} and include an active, direct link to \texttt{https://www.studiva.co.in}.
    \item \textbf{Commercial Restrictions:} Commercial resale, sublicensing, course bundling, or inclusion in paid test series without prior written permission is strictly prohibited.
    \item \textbf{Anti-Tampering Notice:} Removing, altering, or obscuring copyright notices, page watermarks, or website links from this document is strictly prohibited.
\end{itemize}
\vspace{0.6em}
\textbf{Fair Use Disclaimer:}\\
The underlying examination questions originate from the Graduate Aptitude Test in Engineering (GATE), conducted by the Indian Institutes of Technology (IITs) and Indian Institute of Science (IISc). Question texts are reproduced for non-commercial educational study under Fair Use doctrines. No copyright is claimed over the original GATE examination question texts.
\end{tcolorbox}

\vspace{1.2em}
\begin{center}
\small\itshape
Practice Recommendation: Use an opaque sheet of paper or ruler to cover the right-hand column during mock attempts. Tick your chosen answers directly in the check boxes on the left, then slide the cover to check the solution and hint.
\end{center}

\newpage
% ==================== PAGE 2 ONWARDS: QUESTION BANK ====================

\noindent
\begin{tcolorbox}[
    blanker,
    sidebyside,
    sidebyside align=top,
    lefthand width=0.55\textwidth,
    top=4pt, bottom=4pt, left=0pt, right=0pt,
    middle=16pt,
    underlay={
        \draw[dashed, black!80, line width=0.8pt, dash pattern=on 4pt off 3pt]
            ([xshift=0.55\textwidth+8pt]frame.north west) --
            ([xshift=0.55\textwidth+8pt]frame.south west);
    }
]
\textbf{Q.1} \quad {\footnotesize\textsc{[MCQ $\cdot$ 1 Mark(s) $\cdot$ Neg: 0.33]}} \hfill {\footnotesize\textbf{GATE 2018 CS} \quad \tiny\texttt{2018\_CS\_GA\_Q11}}
\vspace{2pt}\hrule height 0.3pt\vspace{4pt}

In an Entity-Relationship (ER) model, suppose $R$ is a many-to-one relationship from entity set $E_1$ to entity set $E_2$. Assume that $E_1$ has total participation in $R$, but $E_2$ has partial participation in $R$. Which one of the following statements is necessarily TRUE?

\par\vspace{6pt}\noindent
\begin{tabular}{@{}p{0.98\linewidth}@{}}
$\square$\ \textbf{(A)} Every entity in $E_1$ is associated with exactly one entity in $E_2$. \\
\addlinespace[3pt]
$\square$\ \textbf{(B)} Some entity in $E_1$ is associated with more than one entity in $E_2$. \\
\addlinespace[3pt]
$\square$\ \textbf{(C)} Every entity in $E_2$ is associated with at least one entity in $E_1$. \\
\addlinespace[3pt]
$\square$\ \textbf{(D)} Every entity in $E_2$ is associated with at most one entity in $E_1$.
\end{tabular}

\tcblower

\textbf{Official Key:} \textbf{\large A} \hfill \textit{\small [Easy]}
\vspace{2pt}\hrule height 0.3pt\vspace{3pt}
\textbf{Topic:} Databases \\
\textbf{Subtopic:} \textit{ER-model}
\vspace{2pt}\hrule height 0.3pt\vspace{3pt}
\textbf{Guide / Hint:} \\
{\small In an ER diagram, a many-to-one relationship from $E_1$ to $E_2$ means each entity in $E_1$ is related to at most one entity in $E_2$. Total participation of $E_1$ requires that every entity in $E_1$ participates in at least one relationship. Combining these, every entity in $E_1$ is associated with exactly one entity in $E_2$.}

\end{tcolorbox}
\vspace{6pt}\hrule height 0.35pt\vspace{6pt}
\noindent
\begin{tcolorbox}[
    blanker,
    sidebyside,
    sidebyside align=top,
    lefthand width=0.55\textwidth,
    top=4pt, bottom=4pt, left=0pt, right=0pt,
    middle=16pt,
    underlay={
        \draw[dashed, black!80, line width=0.8pt, dash pattern=on 4pt off 3pt]
            ([xshift=0.55\textwidth+8pt]frame.north west) --
            ([xshift=0.55\textwidth+8pt]frame.south west);
    }
]
\textbf{Q.2} \quad {\footnotesize\textsc{[MCQ $\cdot$ 1 Mark(s) $\cdot$ Neg: 0.33]}} \hfill {\footnotesize\textbf{GATE 2018 CS} \quad \tiny\texttt{2018\_CS\_GA\_Q12}}
\vspace{2pt}\hrule height 0.3pt\vspace{4pt}

Consider the following two tables and four queries in SQL.
\texttt{Book (isbn, bname)}
\texttt{Stock (isbn, copies)}
\par\medskip
Query 1: \texttt{SELECT B.isbn, S.copies FROM Book B INNER JOIN Stock S ON B.isbn = S.isbn;}
Query 2: \texttt{SELECT B.isbn, S.copies FROM Book B LEFT OUTER JOIN Stock S ON B.isbn = S.isbn;}
Query 3: \texttt{SELECT B.isbn, S.copies FROM Book B RIGHT OUTER JOIN Stock S ON B.isbn = S.isbn;}
Query 4: \texttt{SELECT B.isbn, S.copies FROM Book B FULL OUTER JOIN Stock S ON B.isbn = S.isbn;}
\par\medskip
Which one of the following statements is CORRECT?

\par\vspace{6pt}\noindent
\begin{tabular}{@{}p{0.98\linewidth}@{}}
$\square$\ \textbf{(A)} Query 1 and Query 2 always produce the same number of rows. \\
\addlinespace[3pt]
$\square$\ \textbf{(B)} Query 2 and Query 3 always produce the same number of rows. \\
\addlinespace[3pt]
$\square$\ \textbf{(C)} Query 1 always produces the maximum number of rows among all queries. \\
\addlinespace[3pt]
$\square$\ \textbf{(D)} All rows from Query 1, Query 2, and Query 3 are always present in the result of Query 4.
\end{tabular}

\tcblower

\textbf{Official Key:} \textbf{\large D} \hfill \textit{\small [Easy]}
\vspace{2pt}\hrule height 0.3pt\vspace{3pt}
\textbf{Topic:} Databases \\
\textbf{Subtopic:} \textit{Relational model: relational algebra, tuple calculus, SQL}
\vspace{2pt}\hrule height 0.3pt\vspace{3pt}
\textbf{Guide / Hint:} \\
{\small Query 1 uses natural join (inner join), Query 2 uses left outer join, Query 3 uses right outer join, and Query 4 uses full outer join. A full outer join preserves all tuples from both relation $Book$ and relation $Stock$ (filling missing values with NULLs). Therefore, every tuple output by Query 1, Query 2, and Query 3 is guaranteed to appear in Query 4.}

\end{tcolorbox}
\vspace{6pt}\hrule height 0.35pt\vspace{6pt}
\noindent
\begin{tcolorbox}[
    blanker,
    sidebyside,
    sidebyside align=top,
    lefthand width=0.55\textwidth,
    top=4pt, bottom=4pt, left=0pt, right=0pt,
    middle=16pt,
    underlay={
        \draw[dashed, black!80, line width=0.8pt, dash pattern=on 4pt off 3pt]
            ([xshift=0.55\textwidth+8pt]frame.north west) --
            ([xshift=0.55\textwidth+8pt]frame.south west);
    }
]
\textbf{Q.3} \quad {\footnotesize\textsc{[MCQ $\cdot$ 2 Mark(s) $\cdot$ Neg: 0.66]}} \hfill {\footnotesize\textbf{GATE 2018 CS} \quad \tiny\texttt{2018\_CS\_GA\_Q41}}
\vspace{2pt}\hrule height 0.3pt\vspace{4pt}

Consider the relations $r(A, B)$ and $s(B, C)$, where $s.B$ is a primary key and $r.B$ is a foreign key referencing $s.B$. Consider the query:
\[ Q: \Pi\_{A, B, C}(r \bowtie \sigma\_{B < 5}(s)) \]
Which one of the following relational algebra queries is NOT equivalent to $Q$?

\par\vspace{6pt}\noindent
\begin{tabular}{@{}p{0.98\linewidth}@{}}
$\square$\ \textbf{(A)} $\Pi_{A, B, C}(\sigma_{B < 5}(r) \bowtie s)$ \\
\addlinespace[3pt]
$\square$\ \textbf{(B)} $\sigma_{B < 5}(\Pi_{A, B, C}(r \bowtie s))$ \\
\addlinespace[3pt]
$\square$\ \textbf{(C)} $\Pi_{A, B, C}(r \,\text{⟕}\, \sigma_{B < 5}(s))$ \\
\addlinespace[3pt]
$\square$\ \textbf{(D)} $\Pi_{A, B, C}(\sigma_{B < 5}(r \bowtie s))$
\end{tabular}

\tcblower

\textbf{Official Key:} \textbf{\large C} \hfill \textit{\small [Medium]}
\vspace{2pt}\hrule height 0.3pt\vspace{3pt}
\textbf{Topic:} Databases \\
\textbf{Subtopic:} \textit{Relational model: relational algebra, tuple calculus, SQL}
\vspace{2pt}\hrule height 0.3pt\vspace{3pt}
\textbf{Guide / Hint:} \\
{\small Query $Q$ inner joins $r$ with $\sigma_{B < 5}(s)$, which discards all tuples from $r$ where $B \ge 5$. Option (C) performs a left outer join between $r$ and $\sigma_{B < 5}(s)$, which retains all tuples from $r$ (filling missing $C$ values with NULL for $B \ge 5$), so its result is NOT equivalent to $Q$.}

\end{tcolorbox}
\vspace{6pt}\hrule height 0.35pt\vspace{6pt}
\noindent
\begin{tcolorbox}[
    blanker,
    sidebyside,
    sidebyside align=top,
    lefthand width=0.55\textwidth,
    top=4pt, bottom=4pt, left=0pt, right=0pt,
    middle=16pt,
    underlay={
        \draw[dashed, black!80, line width=0.8pt, dash pattern=on 4pt off 3pt]
            ([xshift=0.55\textwidth+8pt]frame.north west) --
            ([xshift=0.55\textwidth+8pt]frame.south west);
    }
]
\textbf{Q.4} \quad {\footnotesize\textsc{[MCQ $\cdot$ 2 Mark(s) $\cdot$ Neg: 0.66]}} \hfill {\footnotesize\textbf{GATE 2018 CS} \quad \tiny\texttt{2018\_CS\_GA\_Q42}}
\vspace{2pt}\hrule height 0.3pt\vspace{4pt}

Consider the following four relational schemas. For each schema, all non-trivial functional dependencies are listed. The underlined attributes denote candidate keys.
\begin{itemize}[leftmargin=*, nosep]
    \item[Schema I:] $R_1(A, B, C, D)$ with FDs: $A \to B, B \to C, C \to D$
    \item[Schema II:] $R_2(\underline{rollno}, \underline{courseid}, email)$ with FDs: $rollno \to email, email \to rollno$
    \item[Schema III:] $R_3(A, B, C)$ with FDs: $A \to B, B \to C$
    \item[Schema IV:] $R_4(A, B, C)$ with FDs: $AB \to C, C \to A$
\end{itemize}
Which one of the relational schemas is in 3NF but not in BCNF?

\par\vspace{6pt}\noindent
\begin{tabular}{@{}p{0.48\linewidth}@{\hspace{0.04\linewidth}}p{0.48\linewidth}@{}}
$\square$\ \textbf{(A)} Schema I & $\square$\ \textbf{(B)} Schema II \\
\addlinespace[3pt]
$\square$\ \textbf{(C)} Schema III & $\square$\ \textbf{(D)} Schema IV
\end{tabular}

\tcblower

\textbf{Official Key:} \textbf{\large B} \hfill \textit{\small [Medium]}
\vspace{2pt}\hrule height 0.3pt\vspace{3pt}
\textbf{Topic:} Databases \\
\textbf{Subtopic:} \textit{Integrity constraints, normal forms}
\vspace{2pt}\hrule height 0.3pt\vspace{3pt}
\textbf{Guide / Hint:} \\
{\small In Schema II, the functional dependencies are $rollno \to email$ and $email \to rollno$. The candidate keys are $\{rollno, courseid\}$ and $\{email, courseid\}$. In $email \to rollno$, the left side is not a superkey (violating BCNF), but the right side $rollno$ is a prime attribute (satisfying 3NF). Hence Schema II is in 3NF but not in BCNF.}

\end{tcolorbox}
\vspace{6pt}\hrule height 0.35pt\vspace{6pt}
\noindent
\begin{tcolorbox}[
    blanker,
    sidebyside,
    sidebyside align=top,
    lefthand width=0.55\textwidth,
    top=4pt, bottom=4pt, left=0pt, right=0pt,
    middle=16pt,
    underlay={
        \draw[dashed, black!80, line width=0.8pt, dash pattern=on 4pt off 3pt]
            ([xshift=0.55\textwidth+8pt]frame.north west) --
            ([xshift=0.55\textwidth+8pt]frame.south west);
    }
]
\textbf{Q.5} \quad {\footnotesize\textsc{[MCQ $\cdot$ 1 Mark(s) $\cdot$ Neg: 0.33]}} \hfill {\footnotesize\textbf{GATE 2022 CS} \quad \tiny\texttt{2022\_CS\_CS\_Q14}}
\vspace{2pt}\hrule height 0.3pt\vspace{4pt}

In a relational data model, which one of the following statements is TRUE?

\par\vspace{6pt}\noindent
\begin{tabular}{@{}p{0.98\linewidth}@{}}
$\square$\ \textbf{(A)} A relation with only two attributes is always in BCNF. \\
\addlinespace[3pt]
$\square$\ \textbf{(B)} If all attributes of a relation are prime attributes, then the relation is in BCNF. \\
\addlinespace[3pt]
$\square$\ \textbf{(C)} Every relation has at least one non-prime attribute. \\
\addlinespace[3pt]
$\square$\ \textbf{(D)} BCNF decompositions preserve functional dependencies.
\end{tabular}

\tcblower

\textbf{Official Key:} \textbf{\large A} \hfill \textit{\small [Easy]}
\vspace{2pt}\hrule height 0.3pt\vspace{3pt}
\textbf{Topic:} Databases \\
\textbf{Subtopic:} \textit{Integrity constraints, normal forms}
\vspace{2pt}\hrule height 0.3pt\vspace{3pt}
\textbf{Guide / Hint:} \\
{\small For any binary relation $R(A, B)$, the only possible non-trivial functional dependencies are $A \to B$ or $B \to A$. In either case, the left-hand side is a superkey of $R$, satisfying the definition of Boyce-Codd Normal Form (BCNF).}

\end{tcolorbox}
\vspace{6pt}\hrule height 0.35pt\vspace{6pt}
\noindent
\begin{tcolorbox}[
    blanker,
    sidebyside,
    sidebyside align=top,
    lefthand width=0.55\textwidth,
    top=4pt, bottom=4pt, left=0pt, right=0pt,
    middle=16pt,
    underlay={
        \draw[dashed, black!80, line width=0.8pt, dash pattern=on 4pt off 3pt]
            ([xshift=0.55\textwidth+8pt]frame.north west) --
            ([xshift=0.55\textwidth+8pt]frame.south west);
    }
]
\textbf{Q.6} \quad {\footnotesize\textsc{[MSQ $\cdot$ 1 Mark(s) $\cdot$ No Neg]}} \hfill {\footnotesize\textbf{GATE 2022 CS} \quad \tiny\texttt{2022\_CS\_CS\_Q25}}
\vspace{2pt}\hrule height 0.3pt\vspace{4pt}

Consider the following three relations in a relational database:
\[ \text{Employee}(eId, Name), \quad \text{Brand}(bId, bName), \quad \text{Own}(eId, bId) \]
Which of the following relational algebra expressions return the set of $eIds$ who own all the brands?

\par\vspace{6pt}\noindent
\begin{tabular}{@{}p{0.98\linewidth}@{}}
$\square$\ \textbf{(A)} $\Pi_{eId}(\Pi_{eId, bId}(Own) / \Pi_{bId}(Brand))$ \\
\addlinespace[3pt]
$\square$\ \textbf{(B)} $\Pi_{eId}(Own) - \Pi_{eId}((\Pi_{eId}(Own) \times \Pi_{bId}(Brand)) - Own)$ \\
\addlinespace[3pt]
$\square$\ \textbf{(C)} $\Pi_{eId}(\Pi_{eId, bId}(Own) / \Pi_{bId}(Own))$ \\
\addlinespace[3pt]
$\square$\ \textbf{(D)} $\Pi_{eId}((\Pi_{eId}(Own) \times \Pi_{bId}(Brand)) / \Pi_{bId}(Own))$
\end{tabular}

\tcblower

\textbf{Official Key:} \textbf{\large A,B} \hfill \textit{\small [Medium]}
\vspace{2pt}\hrule height 0.3pt\vspace{3pt}
\textbf{Topic:} Databases \\
\textbf{Subtopic:} \textit{Relational model: relational algebra, tuple calculus, SQL}
\vspace{2pt}\hrule height 0.3pt\vspace{3pt}
\textbf{Guide / Hint:} \\
{\small Relational division $\Pi_{eId, bId}(Own) / \Pi_{bId}(Brand)$ directly computes employees who own all brands (A is TRUE). Equivalently, taking all employees, finding those missing at least one brand via cartesian product with Brand followed by set difference with Own, and subtracting them from all employees yields the same set (B is TRUE).}

\end{tcolorbox}
\vspace{6pt}\hrule height 0.35pt\vspace{6pt}
\noindent
\begin{tcolorbox}[
    blanker,
    sidebyside,
    sidebyside align=top,
    lefthand width=0.55\textwidth,
    top=4pt, bottom=4pt, left=0pt, right=0pt,
    middle=16pt,
    underlay={
        \draw[dashed, black!80, line width=0.8pt, dash pattern=on 4pt off 3pt]
            ([xshift=0.55\textwidth+8pt]frame.north west) --
            ([xshift=0.55\textwidth+8pt]frame.south west);
    }
]
\textbf{Q.7} \quad {\footnotesize\textsc{[NAT $\cdot$ 1 Mark(s) $\cdot$ No Neg]}} \hfill {\footnotesize\textbf{GATE 2022 CS} \quad \tiny\texttt{2022\_CS\_CS\_Q31}}
\vspace{2pt}\hrule height 0.3pt\vspace{4pt}

Consider a relation $R(A, B, C, D, E)$ with the following three functional dependencies:
\[ AB \to C; \quad BC \to D; \quad C \to E; \]
The number of superkeys in the relation $R$ is \_\_\_\_\_\_\_\_\_\_\_\. (Answer in integer)

\par\vspace{6pt}\noindent\textbf{Your Answer:} \fbox{\rule{0pt}{1.3ex}\hspace{3cm}}

\tcblower

\textbf{Official Key:} \textbf{\large 8 to 8} \hfill \textit{\small [Easy]}
\vspace{2pt}\hrule height 0.3pt\vspace{3pt}
\textbf{Topic:} Databases \\
\textbf{Subtopic:} \textit{Integrity constraints, normal forms}
\vspace{2pt}\hrule height 0.3pt\vspace{3pt}
\textbf{Guide / Hint:} \\
{\small Computing attribute closures: $(AB)^+ = \{A, B, C, D, E\}$. Attributes $A$ and $B$ never appear on any RHS, so $AB$ is the unique candidate key. Every superkey must contain $AB$, and can include any subset of the remaining 3 attributes $\{C, D, E\}$, giving $2^3 = 8$ superkeys.}

\end{tcolorbox}
\vspace{6pt}\hrule height 0.35pt\vspace{6pt}
\noindent
\begin{tcolorbox}[
    blanker,
    sidebyside,
    sidebyside align=top,
    lefthand width=0.55\textwidth,
    top=4pt, bottom=4pt, left=0pt, right=0pt,
    middle=16pt,
    underlay={
        \draw[dashed, black!80, line width=0.8pt, dash pattern=on 4pt off 3pt]
            ([xshift=0.55\textwidth+8pt]frame.north west) --
            ([xshift=0.55\textwidth+8pt]frame.south west);
    }
]
\textbf{Q.8} \quad {\footnotesize\textsc{[MCQ $\cdot$ 2 Mark(s) $\cdot$ Neg: 0.66]}} \hfill {\footnotesize\textbf{GATE 2022 CS} \quad \tiny\texttt{2022\_CS\_CS\_Q39}}
\vspace{2pt}\hrule height 0.3pt\vspace{4pt}

Let $R_i(z)$ and $W_i(z)$ denote read and write operations on a data element $z$ by a transaction $T_i$, respectively. Consider the schedule $S$ with four transactions:
\[ S: R\_4(x) R\_2(x) R\_3(x) R\_1(y) W\_1(y) W\_2(x) W\_3(y) R\_4(y) \]
Which one of the following serial schedules is conflict equivalent to $S$?

\par\vspace{6pt}\noindent
\begin{tabular}{@{}p{0.48\linewidth}@{\hspace{0.04\linewidth}}p{0.48\linewidth}@{}}
$\square$\ \textbf{(A)} $T_1 \to T_3 \to T_4 \to T_2$ & $\square$\ \textbf{(B)} $T_1 \to T_4 \to T_3 \to T_2$ \\
\addlinespace[3pt]
$\square$\ \textbf{(C)} $T_4 \to T_1 \to T_3 \to T_2$ & $\square$\ \textbf{(D)} $T_3 \to T_1 \to T_4 \to T_2$
\end{tabular}

\tcblower

\textbf{Official Key:} \textbf{\large A} \hfill \textit{\small [Medium]}
\vspace{2pt}\hrule height 0.3pt\vspace{3pt}
\textbf{Topic:} Databases \\
\textbf{Subtopic:} \textit{Transactions and concurrency control}
\vspace{2pt}\hrule height 0.3pt\vspace{3pt}
\textbf{Guide / Hint:} \\
{\small Constructing the serialization graph from conflict operations in $S$: on $y$, $R_1(y), W_1(y) \to W_3(y) \to R_4(y)$ gives edges $T_1 \to T_3 \to T_4$. On $x$, $R_4(x), R_3(x) \to W_2(x)$ gives edges $T_4 \to T_2$ and $T_3 \to T_2$. The unique topological order is $T_1 \to T_3 \to T_4 \to T_2$.}

\end{tcolorbox}
\vspace{6pt}\hrule height 0.35pt\vspace{6pt}
\noindent
\begin{tcolorbox}[
    blanker,
    sidebyside,
    sidebyside align=top,
    lefthand width=0.55\textwidth,
    top=4pt, bottom=4pt, left=0pt, right=0pt,
    middle=16pt,
    underlay={
        \draw[dashed, black!80, line width=0.8pt, dash pattern=on 4pt off 3pt]
            ([xshift=0.55\textwidth+8pt]frame.north west) --
            ([xshift=0.55\textwidth+8pt]frame.south west);
    }
]
\textbf{Q.9} \quad {\footnotesize\textsc{[NAT $\cdot$ 2 Mark(s) $\cdot$ No Neg]}} \hfill {\footnotesize\textbf{GATE 2022 CS} \quad \tiny\texttt{2022\_CS\_CS\_Q56}}
\vspace{2pt}\hrule height 0.3pt\vspace{4pt}

Consider the relational database with the following four schemas and their respective instances:
\[ \text{Student}(sNo, sName, dNo), \quad \text{Dept}(dNo, dName), \quad \text{Course}(cNo, cName, dNo), \quad \text{Register}(sNo, cNo) \]
\begin{lstlisting}[language=SQL]
SELECT * FROM Student AS S WHERE NOT EXISTS
    (SELECT cNo FROM Course WHERE dNo = 'D01'
     EXCEPT
     SELECT cNo FROM Register WHERE sNo = S.sNo)
\end{lstlisting}
The number of rows returned by the above SQL query is \_\_\_\_\_\_\_\_\_\_\_\. (Answer in integer)

\par\vspace{6pt}\noindent\textbf{Your Answer:} \fbox{\rule{0pt}{1.3ex}\hspace{3cm}}

\tcblower

\textbf{Official Key:} \textbf{\large 2 to 2} \hfill \textit{\small [Medium]}
\vspace{2pt}\hrule height 0.3pt\vspace{3pt}
\textbf{Topic:} Databases \\
\textbf{Subtopic:} \textit{Relational model: relational algebra, tuple calculus, SQL}
\vspace{2pt}\hrule height 0.3pt\vspace{3pt}
\textbf{Guide / Hint:} \\
{\small The nested subquery with `EXCEPT` and `NOT EXISTS` expresses relational division: it finds students who have registered for every course offered by department 'D01'. In the database instance, exactly 2 students satisfy this condition.}

\end{tcolorbox}
\vspace{6pt}\hrule height 0.35pt\vspace{6pt}
\noindent
\begin{tcolorbox}[
    blanker,
    sidebyside,
    sidebyside align=top,
    lefthand width=0.55\textwidth,
    top=4pt, bottom=4pt, left=0pt, right=0pt,
    middle=16pt,
    underlay={
        \draw[dashed, black!80, line width=0.8pt, dash pattern=on 4pt off 3pt]
            ([xshift=0.55\textwidth+8pt]frame.north west) --
            ([xshift=0.55\textwidth+8pt]frame.south west);
    }
]
\textbf{Q.10} \quad {\footnotesize\textsc{[MCQ $\cdot$ 1 Mark(s) $\cdot$ Neg: 0.33]}} \hfill {\footnotesize\textbf{GATE 2023 CS} \quad \tiny\texttt{2023\_CS\_CS\_Q16}}
\vspace{2pt}\hrule height 0.3pt\vspace{4pt}

Which one of the options given below refers to the degree (or arity) of a relation in relational database systems?

\par\vspace{6pt}\noindent
\begin{tabular}{@{}p{0.98\linewidth}@{}}
$\square$\ \textbf{(A)} Number of attributes of its relation schema. \\
\addlinespace[3pt]
$\square$\ \textbf{(B)} Number of tuples stored in the relation. \\
\addlinespace[3pt]
$\square$\ \textbf{(C)} Number of entries in the relation. \\
\addlinespace[3pt]
$\square$\ \textbf{(D)} Number of distinct domains of its relation schema.
\end{tabular}

\tcblower

\textbf{Official Key:} \textbf{\large A} \hfill \textit{\small [Easy]}
\vspace{2pt}\hrule height 0.3pt\vspace{3pt}
\textbf{Topic:} Databases \\
\textbf{Subtopic:} \textit{Relational model: relational algebra, tuple calculus, SQL}
\vspace{2pt}\hrule height 0.3pt\vspace{3pt}
\textbf{Guide / Hint:} \\
{\small In relational database theory, the degree (or arity) of a relation is defined as the number of attributes (columns) in its relation schema, whereas the cardinality of a relation is the number of tuples (rows).}

\end{tcolorbox}
\vspace{6pt}\hrule height 0.35pt\vspace{6pt}
\noindent
\begin{tcolorbox}[
    blanker,
    sidebyside,
    sidebyside align=top,
    lefthand width=0.55\textwidth,
    top=4pt, bottom=4pt, left=0pt, right=0pt,
    middle=16pt,
    underlay={
        \draw[dashed, black!80, line width=0.8pt, dash pattern=on 4pt off 3pt]
            ([xshift=0.55\textwidth+8pt]frame.north west) --
            ([xshift=0.55\textwidth+8pt]frame.south west);
    }
]
\textbf{Q.11} \quad {\footnotesize\textsc{[NAT $\cdot$ 2 Mark(s) $\cdot$ No Neg]}} \hfill {\footnotesize\textbf{GATE 2023 CS} \quad \tiny\texttt{2023\_CS\_CS\_Q61}}
\vspace{2pt}\hrule height 0.3pt\vspace{4pt}

Consider the following table named `Student` in a relational database. The primary key of this table is `rollNum`.
\begin{center}
\begin{tabular}{|c|c|c|c|}
\hline
\textbf{rollNum} & \textbf{name} & \textbf{gender} & \textbf{marks} \\
\hline
1 & Naman & M & 62 \\
2 & Aliya & F & 70 \\
3 & Aliya & F & 80 \\
4 & James & M & 82 \\
5 & Swati & F & 65 \\
\hline
\end{tabular}
\end{center}
The SQL query below is executed on this database:
\begin{lstlisting}[language=SQL]
SELECT *
FROM Student
WHERE gender = 'F' AND marks > 65;
\end{lstlisting}
The number of rows returned by the query is \_\_\_\_\_\_\_\_\. (Answer in integer)

\par\vspace{6pt}\noindent\textbf{Your Answer:} \fbox{\rule{0pt}{1.3ex}\hspace{3cm}}

\tcblower

\textbf{Official Key:} \textbf{\large 2 to 2} \hfill \textit{\small [Easy]}
\vspace{2pt}\hrule height 0.3pt\vspace{3pt}
\textbf{Topic:} Databases \\
\textbf{Subtopic:} \textit{Relational model: relational algebra, tuple calculus, SQL}
\vspace{2pt}\hrule height 0.3pt\vspace{3pt}
\textbf{Guide / Hint:} \\
{\small The query evaluates the WHERE clause condition `gender = 'F' AND marks > 65`. Among female students: Aliya (rollNum 2, marks 70) and Aliya (rollNum 3, marks 80) satisfy both conditions, whereas Swati (rollNum 5, marks 65) fails the strict inequality. Thus, exactly 2 rows are returned.}

\end{tcolorbox}
\vspace{6pt}\hrule height 0.35pt\vspace{6pt}
\noindent
\begin{tcolorbox}[
    blanker,
    sidebyside,
    sidebyside align=top,
    lefthand width=0.55\textwidth,
    top=4pt, bottom=4pt, left=0pt, right=0pt,
    middle=16pt,
    underlay={
        \draw[dashed, black!80, line width=0.8pt, dash pattern=on 4pt off 3pt]
            ([xshift=0.55\textwidth+8pt]frame.north west) --
            ([xshift=0.55\textwidth+8pt]frame.south west);
    }
]
\textbf{Q.12} \quad {\footnotesize\textsc{[MCQ $\cdot$ 1 Mark(s) $\cdot$ Neg: 0.33]}} \hfill {\footnotesize\textbf{GATE 2024 CS1} \quad \tiny\texttt{2024\_CS1\_CS\_Q20}}
\vspace{2pt}\hrule height 0.3pt\vspace{4pt}

Let $S$ be the specification: ``Instructors teach courses. Students register for courses. Courses are allocated classrooms. Instructors guide students.'' Which one of the following ER diagrams CORRECTLY represents $S$?

\par\vspace{6pt}\noindent
\begin{tabular}{@{}p{0.48\linewidth}@{\hspace{0.04\linewidth}}p{0.48\linewidth}@{}}
$\square$\ \textbf{(A)} (i) & $\square$\ \textbf{(B)} (ii) \\
\addlinespace[3pt]
$\square$\ \textbf{(C)} (iii) & $\square$\ \textbf{(D)} (iv)
\end{tabular}

\tcblower

\textbf{Official Key:} \textbf{\large D} \hfill \textit{\small [Medium]}
\vspace{2pt}\hrule height 0.3pt\vspace{3pt}
\textbf{Topic:} Databases \\
\textbf{Subtopic:} \textit{ER-model}
\vspace{2pt}\hrule height 0.3pt\vspace{3pt}
\textbf{Guide / Hint:} \\
{\small Specification $S$ describes four distinct binary relationships among entities: Instructors teach courses, Students register for courses, Courses are allocated classrooms, and Instructors guide students. In standard ER modeling, entity sets are represented by rectangles and relationships by diamond diamonds connecting them; diagram (iv) correctly captures all four relationship sets.}

\end{tcolorbox}
\vspace{6pt}\hrule height 0.35pt\vspace{6pt}
\noindent
\begin{tcolorbox}[
    blanker,
    sidebyside,
    sidebyside align=top,
    lefthand width=0.55\textwidth,
    top=4pt, bottom=4pt, left=0pt, right=0pt,
    middle=16pt,
    underlay={
        \draw[dashed, black!80, line width=0.8pt, dash pattern=on 4pt off 3pt]
            ([xshift=0.55\textwidth+8pt]frame.north west) --
            ([xshift=0.55\textwidth+8pt]frame.south west);
    }
]
\textbf{Q.13} \quad {\footnotesize\textsc{[MCQ $\cdot$ 1 Mark(s) $\cdot$ Neg: 0.33]}} \hfill {\footnotesize\textbf{GATE 2024 CS1} \quad \tiny\texttt{2024\_CS1\_CS\_Q21}}
\vspace{2pt}\hrule height 0.3pt\vspace{4pt}

In a $\text{B}^+$ tree, the requirement of at least half-full (50\%) node occupancy is relaxed for which one of the following cases?

\par\vspace{6pt}\noindent
\begin{tabular}{@{}p{0.48\linewidth}@{\hspace{0.04\linewidth}}p{0.48\linewidth}@{}}
$\square$\ \textbf{(A)} Only the root node & $\square$\ \textbf{(B)} All leaf nodes \\
\addlinespace[3pt]
$\square$\ \textbf{(C)} All internal nodes & $\square$\ \textbf{(D)} Only the leftmost leaf node
\end{tabular}

\tcblower

\textbf{Official Key:} \textbf{\large A} \hfill \textit{\small [Easy]}
\vspace{2pt}\hrule height 0.3pt\vspace{3pt}
\textbf{Topic:} Databases \\
\textbf{Subtopic:} \textit{File organization, indexing (e.g., B and B+ trees)}
\vspace{2pt}\hrule height 0.3pt\vspace{3pt}
\textbf{Guide / Hint:} \\
{\small In a $\text{B}^+$ tree, every node except the root node must satisfy the minimum occupancy constraint of being at least half full ($	ext{at least } \lceil p/2 \rceil$ child pointers). The root node is the sole exception, and is permitted to hold as few as two pointers (or one key) when non-leaf, and one key when leaf.}

\end{tcolorbox}
\vspace{6pt}\hrule height 0.35pt\vspace{6pt}
\noindent
\begin{tcolorbox}[
    blanker,
    sidebyside,
    sidebyside align=top,
    lefthand width=0.55\textwidth,
    top=4pt, bottom=4pt, left=0pt, right=0pt,
    middle=16pt,
    underlay={
        \draw[dashed, black!80, line width=0.8pt, dash pattern=on 4pt off 3pt]
            ([xshift=0.55\textwidth+8pt]frame.north west) --
            ([xshift=0.55\textwidth+8pt]frame.south west);
    }
]
\textbf{Q.14} \quad {\footnotesize\textsc{[MSQ $\cdot$ 1 Mark(s) $\cdot$ No Neg]}} \hfill {\footnotesize\textbf{GATE 2024 CS1} \quad \tiny\texttt{2024\_CS1\_CS\_Q22}}
\vspace{2pt}\hrule height 0.3pt\vspace{4pt}

Which of the following statements about a relation $R$ in first normal form (1NF) is/are TRUE?

\par\vspace{6pt}\noindent
\begin{tabular}{@{}p{0.98\linewidth}@{}}
$\square$\ \textbf{(A)} $R$ can have a multi-attribute key \\
\addlinespace[3pt]
$\square$\ \textbf{(B)} $R$ cannot have a foreign key \\
\addlinespace[3pt]
$\square$\ \textbf{(C)} $R$ cannot have a composite attribute \\
\addlinespace[3pt]
$\square$\ \textbf{(D)} $R$ cannot have more than one candidate key
\end{tabular}

\tcblower

\textbf{Official Key:} \textbf{\large A;C} \hfill \textit{\small [Easy]}
\vspace{2pt}\hrule height 0.3pt\vspace{3pt}
\textbf{Topic:} Databases \\
\textbf{Subtopic:} \textit{Integrity constraints, normal forms}
\vspace{2pt}\hrule height 0.3pt\vspace{3pt}
\textbf{Guide / Hint:} \\
{\small A relation is in 1NF if and only if all attribute values are atomic, which explicitly forbids composite or multi-valued attributes (making (C) TRUE). 1NF places no limitation on candidate keys having multiple attributes (multi-attribute keys are fully permitted, making (A) TRUE) or the presence of foreign keys.}

\end{tcolorbox}
\vspace{6pt}\hrule height 0.35pt\vspace{6pt}
\noindent
\begin{tcolorbox}[
    blanker,
    sidebyside,
    sidebyside align=top,
    lefthand width=0.55\textwidth,
    top=4pt, bottom=4pt, left=0pt, right=0pt,
    middle=16pt,
    underlay={
        \draw[dashed, black!80, line width=0.8pt, dash pattern=on 4pt off 3pt]
            ([xshift=0.55\textwidth+8pt]frame.north west) --
            ([xshift=0.55\textwidth+8pt]frame.south west);
    }
]
\textbf{Q.15} \quad {\footnotesize\textsc{[NAT $\cdot$ 1 Mark(s) $\cdot$ No Neg]}} \hfill {\footnotesize\textbf{GATE 2024 CS1} \quad \tiny\texttt{2024\_CS1\_CS\_Q35}}
\vspace{2pt}\hrule height 0.3pt\vspace{4pt}

Consider the following two relations, $R(A,B)$ and $S(A,C)$:
\begin{center}
\begin{tabular}{|c|c|}
\multicolumn{2}{c}{$R$} \\
\hline
$A$ & $B$ \\
\hline
10 & 20 \\
20 & 30 \\
30 & 40 \\
30 & 50 \\
\hline
\end{tabular}
\qquad
\begin{tabular}{|c|c|}
\multicolumn{2}{c}{$S$} \\
\hline
$A$ & $C$ \\
\hline
10 & 90 \\
30 & 45 \\
40 & 80 \\
50 & 95 \\
\hline
\end{tabular}
\end{center}
The total number of tuples obtained by evaluating the following expression
\[ \sigma\_{B < C}(R \bowtie\_{R.A = S.A} S) \]
is \_\_\_\_\_\_\_\_\_\. (Answer in integer)

\par\vspace{6pt}\noindent\textbf{Your Answer:} \fbox{\rule{0pt}{1.3ex}\hspace{3cm}}

\tcblower

\textbf{Official Key:} \textbf{\large 2 to 2} \hfill \textit{\small [Easy]}
\vspace{2pt}\hrule height 0.3pt\vspace{3pt}
\textbf{Topic:} Databases \\
\textbf{Subtopic:} \textit{Relational model: relational algebra, tuple calculus, SQL}
\vspace{2pt}\hrule height 0.3pt\vspace{3pt}
\textbf{Guide / Hint:} \\
{\small Joining $R(A,B)$ and $S(A,C)$ on $R.A = S.A$ yields: $(10, 20, 90)$ for $A=10$, and $(30, 40, 45)$, $(30, 50, 45)$ for $A=30$. Applying selection $\sigma_{B < C}$: $20 < 90$ is true (tuple $(10, 20, 90)$), $40 < 45$ is true (tuple $(30, 40, 45)$), and $50 < 45$ is false. Thus, exactly 2 tuples satisfy the condition.}

\end{tcolorbox}
\vspace{6pt}\hrule height 0.35pt\vspace{6pt}
\noindent
\begin{tcolorbox}[
    blanker,
    sidebyside,
    sidebyside align=top,
    lefthand width=0.55\textwidth,
    top=4pt, bottom=4pt, left=0pt, right=0pt,
    middle=16pt,
    underlay={
        \draw[dashed, black!80, line width=0.8pt, dash pattern=on 4pt off 3pt]
            ([xshift=0.55\textwidth+8pt]frame.north west) --
            ([xshift=0.55\textwidth+8pt]frame.south west);
    }
]
\textbf{Q.16} \quad {\footnotesize\textsc{[MSQ $\cdot$ 2 Mark(s) $\cdot$ No Neg]}} \hfill {\footnotesize\textbf{GATE 2024 CS1} \quad \tiny\texttt{2024\_CS1\_CS\_Q44}}
\vspace{2pt}\hrule height 0.3pt\vspace{4pt}

The symbol $\to$ indicates functional dependency in the context of a relational database. Which of the following options is/are TRUE?

\par\vspace{6pt}\noindent
\begin{tabular}{@{}p{0.98\linewidth}@{}}
$\square$\ \textbf{(A)} $(X, Y) \to (Z, W)$ implies $X \to (Z, W)$ \\
\addlinespace[3pt]
$\square$\ \textbf{(B)} $(X, Y) \to (Z, W)$ implies $(X, Y) \to Z$ \\
\addlinespace[3pt]
$\square$\ \textbf{(C)} $((X, Y) \to Z \text{ and } W \to Y)$ implies $(X, W) \to Z$ \\
\addlinespace[3pt]
$\square$\ \textbf{(D)} $(X \to Y \text{ and } Y \to Z)$ implies $X \to Z$
\end{tabular}

\tcblower

\textbf{Official Key:} \textbf{\large B;C;D} \hfill \textit{\small [Medium]}
\vspace{2pt}\hrule height 0.3pt\vspace{3pt}
\textbf{Topic:} Databases \\
\textbf{Subtopic:} \textit{Integrity constraints, normal forms}
\vspace{2pt}\hrule height 0.3pt\vspace{3pt}
\textbf{Guide / Hint:} \\
{\small By Armstrong's decomposition axiom, $(X, Y) \to (Z, W)$ implies $(X, Y) \to Z$ (B is TRUE). By Armstrong's transitivity axiom, $X \to Y$ and $Y \to Z$ implies $X \to Z$ (D is TRUE). By pseudotransitivity, $(X, Y) \to Z$ and $W \to Y$ implies $(X, W) \to Z$ (C is TRUE). Decomposing the left-hand side determinant is invalid in general (A is FALSE).}

\end{tcolorbox}
\vspace{6pt}\hrule height 0.35pt\vspace{6pt}
\noindent
\begin{tcolorbox}[
    blanker,
    sidebyside,
    sidebyside align=top,
    lefthand width=0.55\textwidth,
    top=4pt, bottom=4pt, left=0pt, right=0pt,
    middle=16pt,
    underlay={
        \draw[dashed, black!80, line width=0.8pt, dash pattern=on 4pt off 3pt]
            ([xshift=0.55\textwidth+8pt]frame.north west) --
            ([xshift=0.55\textwidth+8pt]frame.south west);
    }
]
\textbf{Q.17} \quad {\footnotesize\textsc{[MSQ $\cdot$ 2 Mark(s) $\cdot$ No Neg]}} \hfill {\footnotesize\textbf{GATE 2024 CS1} \quad \tiny\texttt{2024\_CS1\_CS\_Q46}}
\vspace{2pt}\hrule height 0.3pt\vspace{4pt}

Consider the following read-write schedule $S$ over three transactions $T_1, T_2$, and $T_3$, where the subscripts in the schedule indicate transaction IDs:
\[ S: r\_1(z); w\_1(z); r\_2(x); r\_3(y); w\_3(y); r\_2(y); w\_2(x); w\_2(y); \]
Which of the following transaction schedules is/are conflict equivalent to $S$?

\par\vspace{6pt}\noindent
\begin{tabular}{@{}p{0.48\linewidth}@{\hspace{0.04\linewidth}}p{0.48\linewidth}@{}}
$\square$\ \textbf{(A)} $T_1 T_2 T_3$ & $\square$\ \textbf{(B)} $T_1 T_3 T_2$ \\
\addlinespace[3pt]
$\square$\ \textbf{(C)} $T_3 T_2 T_1$ & $\square$\ \textbf{(D)} $T_3 T_1 T_2$
\end{tabular}

\tcblower

\textbf{Official Key:} \textbf{\large B;C;D} \hfill \textit{\small [Medium]}
\vspace{2pt}\hrule height 0.3pt\vspace{3pt}
\textbf{Topic:} Databases \\
\textbf{Subtopic:} \textit{Transactions and concurrency control}
\vspace{2pt}\hrule height 0.3pt\vspace{3pt}
\textbf{Guide / Hint:} \\
{\small Examining data item conflicts: on item $y$, $w_3(y)$ in $T_3$ precedes both $r_2(y)$ and $w_2(y)$ in $T_2$, creating a precedence conflict edge $T_3 \to T_2$. Transaction $T_1$ only accesses $z$ and has no conflicting operations with $T_2$ or $T_3$. Any serial schedule preserving $T_3$ before $T_2$ is conflict equivalent to $S$, namely $T_1 T_3 T_2$, $T_3 T_1 T_2$, and $T_3 T_2 T_1$.}

\end{tcolorbox}
\vspace{6pt}\hrule height 0.35pt\vspace{6pt}
\noindent
\begin{tcolorbox}[
    blanker,
    sidebyside,
    sidebyside align=top,
    lefthand width=0.55\textwidth,
    top=4pt, bottom=4pt, left=0pt, right=0pt,
    middle=16pt,
    underlay={
        \draw[dashed, black!80, line width=0.8pt, dash pattern=on 4pt off 3pt]
            ([xshift=0.55\textwidth+8pt]frame.north west) --
            ([xshift=0.55\textwidth+8pt]frame.south west);
    }
]
\textbf{Q.18} \quad {\footnotesize\textsc{[MCQ $\cdot$ 1 Mark(s) $\cdot$ Neg: 0.33]}} \hfill {\footnotesize\textbf{GATE 2024 CS2} \quad \tiny\texttt{2024\_CS2\_CS\_Q19}}
\vspace{2pt}\hrule height 0.3pt\vspace{4pt}

Once the DBMS informs the user that a transaction has been successfully completed, its effect should persist even if the system crashes before all its changes are reflected on disk. This property is called:

\par\vspace{6pt}\noindent
\begin{tabular}{@{}p{0.48\linewidth}@{\hspace{0.04\linewidth}}p{0.48\linewidth}@{}}
$\square$\ \textbf{(A)} durability & $\square$\ \textbf{(B)} atomicity \\
\addlinespace[3pt]
$\square$\ \textbf{(C)} consistency & $\square$\ \textbf{(D)} isolation
\end{tabular}

\tcblower

\textbf{Official Key:} \textbf{\large A} \hfill \textit{\small [Easy]}
\vspace{2pt}\hrule height 0.3pt\vspace{3pt}
\textbf{Topic:} Databases \\
\textbf{Subtopic:} \textit{Transactions and concurrency control}
\vspace{2pt}\hrule height 0.3pt\vspace{3pt}
\textbf{Guide / Hint:} \\
{\small The Durability property of the ACID paradigm guarantees that once a transaction has committed successfully, its updates and effects survive permanently in non-volatile storage, even in the event of a subsequent system crash.}

\end{tcolorbox}
\vspace{6pt}\hrule height 0.35pt\vspace{6pt}
\noindent
\begin{tcolorbox}[
    blanker,
    sidebyside,
    sidebyside align=top,
    lefthand width=0.55\textwidth,
    top=4pt, bottom=4pt, left=0pt, right=0pt,
    middle=16pt,
    underlay={
        \draw[dashed, black!80, line width=0.8pt, dash pattern=on 4pt off 3pt]
            ([xshift=0.55\textwidth+8pt]frame.north west) --
            ([xshift=0.55\textwidth+8pt]frame.south west);
    }
]
\textbf{Q.19} \quad {\footnotesize\textsc{[MCQ $\cdot$ 1 Mark(s) $\cdot$ Neg: 0.33]}} \hfill {\footnotesize\textbf{GATE 2024 CS2} \quad \tiny\texttt{2024\_CS2\_CS\_Q20}}
\vspace{2pt}\hrule height 0.3pt\vspace{4pt}

In the context of owner and weak entity sets in the ER (Entity-Relationship) data model, which one of the following statements is TRUE?

\par\vspace{6pt}\noindent
\begin{tabular}{@{}p{0.98\linewidth}@{}}
$\square$\ \textbf{(A)} The weak entity set MUST have total participation in the identifying relationship \\
\addlinespace[3pt]
$\square$\ \textbf{(B)} The owner entity set MUST have total participation in the identifying relationship \\
\addlinespace[3pt]
$\square$\ \textbf{(C)} Both weak and owner entity sets MUST have total participation in the identifying relationship \\
\addlinespace[3pt]
$\square$\ \textbf{(D)} Neither weak entity set nor owner entity set MUST have total participation in the identifying relationship
\end{tabular}

\tcblower

\textbf{Official Key:} \textbf{\large A} \hfill \textit{\small [Easy]}
\vspace{2pt}\hrule height 0.3pt\vspace{3pt}
\textbf{Topic:} Databases \\
\textbf{Subtopic:} \textit{ER-model}
\vspace{2pt}\hrule height 0.3pt\vspace{3pt}
\textbf{Guide / Hint:} \\
{\small A weak entity cannot exist without its identifying (owner) entity. Therefore, every instance of the weak entity set must be associated with an owner entity instance via the identifying relationship, mandating total participation of the weak entity set.}

\end{tcolorbox}
\vspace{6pt}\hrule height 0.35pt\vspace{6pt}
\noindent
\begin{tcolorbox}[
    blanker,
    sidebyside,
    sidebyside align=top,
    lefthand width=0.55\textwidth,
    top=4pt, bottom=4pt, left=0pt, right=0pt,
    middle=16pt,
    underlay={
        \draw[dashed, black!80, line width=0.8pt, dash pattern=on 4pt off 3pt]
            ([xshift=0.55\textwidth+8pt]frame.north west) --
            ([xshift=0.55\textwidth+8pt]frame.south west);
    }
]
\textbf{Q.20} \quad {\footnotesize\textsc{[MSQ $\cdot$ 1 Mark(s) $\cdot$ No Neg]}} \hfill {\footnotesize\textbf{GATE 2024 CS2} \quad \tiny\texttt{2024\_CS2\_CS\_Q26}}
\vspace{2pt}\hrule height 0.3pt\vspace{4pt}

Which of the following file organizations is/are I/O efficient for the scan operation in DBMS?

\par\vspace{6pt}\noindent
\begin{tabular}{@{}p{0.48\linewidth}@{\hspace{0.04\linewidth}}p{0.48\linewidth}@{}}
$\square$\ \textbf{(A)} Sorted & $\square$\ \textbf{(B)} Heap \\
\addlinespace[3pt]
$\square$\ \textbf{(C)} Unclustered tree index & $\square$\ \textbf{(D)} Unclustered hash index
\end{tabular}

\tcblower

\textbf{Official Key:} \textbf{\large A;B} \hfill \textit{\small [Easy]}
\vspace{2pt}\hrule height 0.3pt\vspace{3pt}
\textbf{Topic:} Databases \\
\textbf{Subtopic:} \textit{File organization, indexing (e.g., B and B+ trees)}
\vspace{2pt}\hrule height 0.3pt\vspace{3pt}
\textbf{Guide / Hint:} \\
{\small A full table scan retrieves all pages sequentially. Sorted and Heap files store data records densely packed across consecutive data pages, requiring reading each page exactly once. Unclustered indexes store data out of order, incurring multiple random I/Os per record.}

\end{tcolorbox}
\vspace{6pt}\hrule height 0.35pt\vspace{6pt}
\noindent
\begin{tcolorbox}[
    blanker,
    sidebyside,
    sidebyside align=top,
    lefthand width=0.55\textwidth,
    top=4pt, bottom=4pt, left=0pt, right=0pt,
    middle=16pt,
    underlay={
        \draw[dashed, black!80, line width=0.8pt, dash pattern=on 4pt off 3pt]
            ([xshift=0.55\textwidth+8pt]frame.north west) --
            ([xshift=0.55\textwidth+8pt]frame.south west);
    }
]
\textbf{Q.21} \quad {\footnotesize\textsc{[MSQ $\cdot$ 1 Mark(s) $\cdot$ No Neg]}} \hfill {\footnotesize\textbf{GATE 2024 CS2} \quad \tiny\texttt{2024\_CS2\_CS\_Q27}}
\vspace{2pt}\hrule height 0.3pt\vspace{4pt}

Which of the following statements about the Two Phase Locking (2PL) protocol is/are TRUE?

\par\vspace{6pt}\noindent
\begin{tabular}{@{}p{0.98\linewidth}@{}}
$\square$\ \textbf{(A)} 2PL permits only serializable schedules \\
\addlinespace[3pt]
$\square$\ \textbf{(B)} With 2PL, a transaction always locks the data item being read or written just before every operation \\
\addlinespace[3pt]
$\square$\ \textbf{(C)} With 2PL, once a lock is released on any data item inside a transaction, no more locks on any data item can be acquired \\
\addlinespace[3pt]
$\square$\ \textbf{(D)} A deadlock is possible with 2PL
\end{tabular}

\tcblower

\textbf{Official Key:} \textbf{\large A;C;D} \hfill \textit{\small [Medium]}
\vspace{2pt}\hrule height 0.3pt\vspace{3pt}
\textbf{Topic:} Databases \\
\textbf{Subtopic:} \textit{Transactions and concurrency control}
\vspace{2pt}\hrule height 0.3pt\vspace{3pt}
\textbf{Guide / Hint:} \\
{\small Two-Phase Locking (2PL) guarantees conflict serializability (A). In 2PL, once a lock is released (shrinking phase), no new locks can be acquired (C). Basic 2PL is susceptible to deadlocks due to circular wait conditions (D). Statement B describes conservative/static locking, not general 2PL.}

\end{tcolorbox}
\vspace{6pt}\hrule height 0.35pt\vspace{6pt}
\noindent
\begin{tcolorbox}[
    blanker,
    sidebyside,
    sidebyside align=top,
    lefthand width=0.55\textwidth,
    top=4pt, bottom=4pt, left=0pt, right=0pt,
    middle=16pt,
    underlay={
        \draw[dashed, black!80, line width=0.8pt, dash pattern=on 4pt off 3pt]
            ([xshift=0.55\textwidth+8pt]frame.north west) --
            ([xshift=0.55\textwidth+8pt]frame.south west);
    }
]
\textbf{Q.22} \quad {\footnotesize\textsc{[MCQ $\cdot$ 2 Mark(s) $\cdot$ Neg: 0.66]}} \hfill {\footnotesize\textbf{GATE 2024 CS2} \quad \tiny\texttt{2024\_CS2\_CS\_Q45}}
\vspace{2pt}\hrule height 0.3pt\vspace{4pt}

The relation schema, Person(pid, city), describes the city of residence for every person uniquely identified by pid. The following relational algebra operators are available: selection, projection, cross product, and rename.
To find the list of cities where at least 3 persons reside, what is the minimum number of cross product operations required?

\par\vspace{6pt}\noindent
\begin{tabular}{@{}p{0.48\linewidth}@{\hspace{0.04\linewidth}}p{0.48\linewidth}@{}}
$\square$\ \textbf{(A)} 1 & $\square$\ \textbf{(B)} 2 \\
\addlinespace[3pt]
$\square$\ \textbf{(C)} 3 & $\square$\ \textbf{(D)} 4
\end{tabular}

\tcblower

\textbf{Official Key:} \textbf{\large B} \hfill \textit{\small [Medium]}
\vspace{2pt}\hrule height 0.3pt\vspace{3pt}
\textbf{Topic:} Databases \\
\textbf{Subtopic:} \textit{Relational model: relational algebra, tuple calculus, SQL}
\vspace{2pt}\hrule height 0.3pt\vspace{3pt}
\textbf{Guide / Hint:} \\
{\small To find cities with at least 3 distinct persons residing in them using selection, projection, cross-product, and rename, we need 3 copies of the Person table ($P_1, P_2, P_3$). Combining 3 relations requires 2 cross-product operations: $(P_1 \times P_2) \times P_3$.}

\end{tcolorbox}
\vspace{6pt}\hrule height 0.35pt\vspace{6pt}
\noindent
\begin{tcolorbox}[
    blanker,
    sidebyside,
    sidebyside align=top,
    lefthand width=0.55\textwidth,
    top=4pt, bottom=4pt, left=0pt, right=0pt,
    middle=16pt,
    underlay={
        \draw[dashed, black!80, line width=0.8pt, dash pattern=on 4pt off 3pt]
            ([xshift=0.55\textwidth+8pt]frame.north west) --
            ([xshift=0.55\textwidth+8pt]frame.south west);
    }
]
\textbf{Q.23} \quad {\footnotesize\textsc{[NAT $\cdot$ 2 Mark(s) $\cdot$ No Neg]}} \hfill {\footnotesize\textbf{GATE 2024 CS2} \quad \tiny\texttt{2024\_CS2\_CS\_Q56}}
\vspace{2pt}\hrule height 0.3pt\vspace{4pt}

A functional dependency $F: X \to Y$ is termed as a useful functional dependency if and only if it satisfies all the following three conditions:
\begin{itemize}
\item $X$ is not the empty set.
\item $Y$ is not the empty set.
\item Intersection of $X$ and $Y$ is the empty set.
\end{itemize}
For a relation $R$ with 4 attributes, the total number of possible useful functional dependencies is \_\_\_\_\_\_\. (Answer in integer)

\par\vspace{6pt}\noindent\textbf{Your Answer:} \fbox{\rule{0pt}{1.3ex}\hspace{3cm}}

\tcblower

\textbf{Official Key:} \textbf{\large 50 to 50} \hfill \textit{\small [Medium]}
\vspace{2pt}\hrule height 0.3pt\vspace{3pt}
\textbf{Topic:} Databases \\
\textbf{Subtopic:} \textit{Integrity constraints, normal forms}
\vspace{2pt}\hrule height 0.3pt\vspace{3pt}
\textbf{Guide / Hint:} \\
{\small For each attribute $A$ among the 4 attributes, there are 3 mutually exclusive choices: $A \in X$, $A \in Y$, or $A \notin (X \cup Y)$, giving $3^4 = 81$ pairs of disjoint sets $(X, Y)$. We subtract cases where $X = \emptyset$ ($2^4 = 16$), where $Y = \emptyset$ ($2^4 = 16$), and add back the double-subtracted case where $X = Y = \emptyset$ ($1$). Total useful FDs $= 81 - 16 - 16 + 1 = 50$.}

\end{tcolorbox}
\vspace{6pt}\hrule height 0.35pt\vspace{6pt}
\noindent
\begin{tcolorbox}[
    blanker,
    sidebyside,
    sidebyside align=top,
    lefthand width=0.55\textwidth,
    top=4pt, bottom=4pt, left=0pt, right=0pt,
    middle=16pt,
    underlay={
        \draw[dashed, black!80, line width=0.8pt, dash pattern=on 4pt off 3pt]
            ([xshift=0.55\textwidth+8pt]frame.north west) --
            ([xshift=0.55\textwidth+8pt]frame.south west);
    }
]
\textbf{Q.24} \quad {\footnotesize\textsc{[MCQ $\cdot$ 1 Mark(s) $\cdot$ Neg: 0.33]}} \hfill {\footnotesize\textbf{GATE 2025 CS1} \quad \tiny\texttt{2025\_CS1\_CS\_Q15}}
\vspace{2pt}\hrule height 0.3pt\vspace{4pt}

A schedule of three database transactions $T_1, T_2$, and $T_3$ is shown. $R_i(A)$ and $W_i(A)$ denote read and write of data item $A$ by transaction $T_i$, $i= 1,2,3$. The transaction $T_1$ aborts at the end. Which other transaction(s) will be required to be rolled back?
\[ R\_1(X) \; W\_1(Y) \; R\_2(X) \; R\_2(Y) \; R\_3(Y) \; ABORT(T\_1) \]

\par\vspace{6pt}\noindent
\begin{tabular}{@{}p{0.48\linewidth}@{\hspace{0.04\linewidth}}p{0.48\linewidth}@{}}
$\square$\ \textbf{(A)} Only $T_2$ & $\square$\ \textbf{(B)} Only $T_3$ \\
\addlinespace[3pt]
$\square$\ \textbf{(C)} Both $T_2$ and $T_3$ & $\square$\ \textbf{(D)} Neither $T_2$ nor $T_3$
\end{tabular}

\tcblower

\textbf{Official Key:} \textbf{\large C} \hfill \textit{\small [Medium]}
\vspace{2pt}\hrule height 0.3pt\vspace{3pt}
\textbf{Topic:} Databases \\
\textbf{Subtopic:} \textit{Transactions and concurrency control}
\vspace{2pt}\hrule height 0.3pt\vspace{3pt}
\textbf{Guide / Hint:} \\
{\small Transaction $T_1$ writes to item $Y$ ($W_1(Y)$). Subsequently, both $T_2$ and $T_3$ read uncommitted data item $Y$ ($R_2(Y)$ and $R_3(Y)$) written by $T_1$ before $T_1$ commits. When $T_1$ aborts, a cascading rollback is triggered, requiring both $T_2$ and $T_3$ to be rolled back.}

\end{tcolorbox}
\vspace{6pt}\hrule height 0.35pt\vspace{6pt}
\noindent
\begin{tcolorbox}[
    blanker,
    sidebyside,
    sidebyside align=top,
    lefthand width=0.55\textwidth,
    top=4pt, bottom=4pt, left=0pt, right=0pt,
    middle=16pt,
    underlay={
        \draw[dashed, black!80, line width=0.8pt, dash pattern=on 4pt off 3pt]
            ([xshift=0.55\textwidth+8pt]frame.north west) --
            ([xshift=0.55\textwidth+8pt]frame.south west);
    }
]
\textbf{Q.25} \quad {\footnotesize\textsc{[MSQ $\cdot$ 1 Mark(s) $\cdot$ No Neg]}} \hfill {\footnotesize\textbf{GATE 2025 CS1} \quad \tiny\texttt{2025\_CS1\_CS\_Q21}}
\vspace{2pt}\hrule height 0.3pt\vspace{4pt}

Consider a $B^+$ tree with 5 nodes, in which a node can store at most 3 key values. The value 23 is now inserted in the $B^+$ tree. Which of the following option(s) is/are CORRECT?

\par\vspace{6pt}\noindent
\begin{tabular}{@{}p{0.98\linewidth}@{}}
$\square$\ \textbf{(A)} None of the nodes will split. \\
\addlinespace[3pt]
$\square$\ \textbf{(B)} At least one node will split and redistribute. \\
\addlinespace[3pt]
$\square$\ \textbf{(C)} The total number of nodes will remain same. \\
\addlinespace[3pt]
$\square$\ \textbf{(D)} The height of the tree will increase.
\end{tabular}

\tcblower

\textbf{Official Key:} \textbf{\large B;D} \hfill \textit{\small [Medium]}
\vspace{2pt}\hrule height 0.3pt\vspace{3pt}
\textbf{Topic:} Databases \\
\textbf{Subtopic:} \textit{File organization, indexing (e.g., B and B+ trees)}
\vspace{2pt}\hrule height 0.3pt\vspace{3pt}
\textbf{Guide / Hint:} \\
{\small With maximum key capacity of 3 keys per node, the targeted leaf node containing keys $[20, 21, 22]$ is already full. Inserting 23 causes this leaf node to overflow and split, propagating a split upward to the already-saturated root, which creates a new root node and increases tree height.}

\end{tcolorbox}
\vspace{6pt}\hrule height 0.35pt\vspace{6pt}
\noindent
\begin{tcolorbox}[
    blanker,
    sidebyside,
    sidebyside align=top,
    lefthand width=0.55\textwidth,
    top=4pt, bottom=4pt, left=0pt, right=0pt,
    middle=16pt,
    underlay={
        \draw[dashed, black!80, line width=0.8pt, dash pattern=on 4pt off 3pt]
            ([xshift=0.55\textwidth+8pt]frame.north west) --
            ([xshift=0.55\textwidth+8pt]frame.south west);
    }
]
\textbf{Q.26} \quad {\footnotesize\textsc{[MCQ $\cdot$ 2 Mark(s) $\cdot$ Neg: 0.66]}} \hfill {\footnotesize\textbf{GATE 2025 CS1} \quad \tiny\texttt{2025\_CS1\_CS\_Q39}}
\vspace{2pt}\hrule height 0.3pt\vspace{4pt}

Consider two relations describing $teams$ and $players$ in a sports league:
\begin{itemize}
\item $teams(tid, tname)$: $tid, tname$ are team-id and team-name, respectively
\item $players(pid, pname, tid)$: $pid, pname$, and $tid$ denote player-id, player-name and the team-id of the player, respectively
\end{itemize}
Which ONE of the following tuple relational calculus queries returns the name of the players who play for the team having $tname$ as 'MI'?

\par\vspace{6pt}\noindent
\begin{tabular}{@{}p{0.98\linewidth}@{}}
$\square$\ \textbf{(A)} $\{ p.pname \mid p \in players \land \exists t (t \in teams \land p.tid = t.tid \land t.tname = \text{'MI'}) \}$ \\
\addlinespace[3pt]
$\square$\ \textbf{(B)} $\{ p.pname \mid p \in teams \land \exists t (t \in players \land p.tid = t.tid \land t.tname = \text{'MI'}) \}$ \\
\addlinespace[3pt]
$\square$\ \textbf{(C)} $\{ p.pname \mid p \in players \land \exists t (t \in teams \land t.tname = \text{'MI'}) \}$ \\
\addlinespace[3pt]
$\square$\ \textbf{(D)} $\{ p.pname \mid p \in teams \land \exists t (t \in players \land t.tname = \text{'MI'}) \}$
\end{tabular}

\tcblower

\textbf{Official Key:} \textbf{\large A} \hfill \textit{\small [Medium]}
\vspace{2pt}\hrule height 0.3pt\vspace{3pt}
\textbf{Topic:} Databases \\
\textbf{Subtopic:} \textit{Relational model: relational algebra, tuple calculus, SQL}
\vspace{2pt}\hrule height 0.3pt\vspace{3pt}
\textbf{Guide / Hint:} \\
{\small The query must select `p.pname` from the `players` relation ($p \in players$) and join with a tuple $t \in teams$ on the foreign key condition $p.tid = t.tid$, while filtering for $t.tname = \text{'MI'}$. Option A accurately expresses this quantified condition in Tuple Relational Calculus.}

\end{tcolorbox}
\vspace{6pt}\hrule height 0.35pt\vspace{6pt}
\noindent
\begin{tcolorbox}[
    blanker,
    sidebyside,
    sidebyside align=top,
    lefthand width=0.55\textwidth,
    top=4pt, bottom=4pt, left=0pt, right=0pt,
    middle=16pt,
    underlay={
        \draw[dashed, black!80, line width=0.8pt, dash pattern=on 4pt off 3pt]
            ([xshift=0.55\textwidth+8pt]frame.north west) --
            ([xshift=0.55\textwidth+8pt]frame.south west);
    }
]
\textbf{Q.27} \quad {\footnotesize\textsc{[MSQ $\cdot$ 2 Mark(s) $\cdot$ No Neg]}} \hfill {\footnotesize\textbf{GATE 2025 CS1} \quad \tiny\texttt{2025\_CS1\_CS\_Q47}}
\vspace{2pt}\hrule height 0.3pt\vspace{4pt}

Consider a relational schema $team(name, city, owner)$, with functional dependencies $\{name \to city, name \to owner\}$. The relation $team$ is decomposed into two relations, $t_1(name, city)$ and $t_2(name, owner)$. Which of the following statement(s) is/are TRUE?

\par\vspace{6pt}\noindent
\begin{tabular}{@{}p{0.98\linewidth}@{}}
$\square$\ \textbf{(A)} The relation $team$ is NOT in BCNF. \\
\addlinespace[3pt]
$\square$\ \textbf{(B)} The relations $t_1$ and $t_2$ are in BCNF. \\
\addlinespace[3pt]
$\square$\ \textbf{(C)} The decomposition constitutes a lossless join. \\
\addlinespace[3pt]
$\square$\ \textbf{(D)} The relation $team$ is NOT in 3NF.
\end{tabular}

\tcblower

\textbf{Official Key:} \textbf{\large B;C} \hfill \textit{\small [Medium]}
\vspace{2pt}\hrule height 0.3pt\vspace{3pt}
\textbf{Topic:} Databases \\
\textbf{Subtopic:} \textit{Integrity constraints, normal forms}
\vspace{2pt}\hrule height 0.3pt\vspace{3pt}
\textbf{Guide / Hint:} \\
{\small In $team$, $name$ is a candidate key since $name^+ = \{name, city, owner\}$. Since every FD has a superkey on the LHS, $team$ was already in BCNF and 3NF (A and D are false). In $t_1$ and $t_2$, $name$ is the primary key, so both are in BCNF (B is true). The common attribute $R_1 \cap R_2 = \{name\}$ is a key of both $t_1$ and $t_2$, ensuring a lossless join (C is true).}

\end{tcolorbox}
\vspace{6pt}\hrule height 0.35pt\vspace{6pt}
\noindent
\begin{tcolorbox}[
    blanker,
    sidebyside,
    sidebyside align=top,
    lefthand width=0.55\textwidth,
    top=4pt, bottom=4pt, left=0pt, right=0pt,
    middle=16pt,
    underlay={
        \draw[dashed, black!80, line width=0.8pt, dash pattern=on 4pt off 3pt]
            ([xshift=0.55\textwidth+8pt]frame.north west) --
            ([xshift=0.55\textwidth+8pt]frame.south west);
    }
]
\textbf{Q.28} \quad {\footnotesize\textsc{[MSQ $\cdot$ 1 Mark(s) $\cdot$ No Neg]}} \hfill {\footnotesize\textbf{GATE 2025 CS2} \quad \tiny\texttt{2025\_CS2\_CS\_Q27}}
\vspace{2pt}\hrule height 0.3pt\vspace{4pt}

An audit of a banking transactions system has found that on an earlier occasion, two joint holders of account $A$ attempted simultaneous transfers of Rs. 10000 each from account $A$ to account $B$. Both transactions read the same value, Rs. 11000, as the initial balance in $A$ and were allowed to go through. $B$ was credited Rs. 10000 twice. $A$ was debited only once and ended up with a balance of Rs. 1000. Which of the following properties is/are certain to have been violated by the system?

\par\vspace{6pt}\noindent
\begin{tabular}{@{}p{0.48\linewidth}@{\hspace{0.04\linewidth}}p{0.48\linewidth}@{}}
$\square$\ \textbf{(A)} Atomicity & $\square$\ \textbf{(B)} Consistency \\
\addlinespace[3pt]
$\square$\ \textbf{(C)} Isolation & $\square$\ \textbf{(D)} Durability
\end{tabular}

\tcblower

\textbf{Official Key:} \textbf{\large B;C} \hfill \textit{\small [Medium]}
\vspace{2pt}\hrule height 0.3pt\vspace{3pt}
\textbf{Topic:} Databases \\
\textbf{Subtopic:} \textit{Transactions and concurrency control: ACID properties}
\vspace{2pt}\hrule height 0.3pt\vspace{3pt}
\textbf{Guide / Hint:} \\
{\small Both transactions read the same uncommitted state and overwrote each other (lost update anomaly), violating Isolation. The final balance violates the database invariant (starting 11000 debited $2 \times 10000$ cannot yield 1000), violating Consistency.}

\end{tcolorbox}
\vspace{6pt}\hrule height 0.35pt\vspace{6pt}
\noindent
\begin{tcolorbox}[
    blanker,
    sidebyside,
    sidebyside align=top,
    lefthand width=0.55\textwidth,
    top=4pt, bottom=4pt, left=0pt, right=0pt,
    middle=16pt,
    underlay={
        \draw[dashed, black!80, line width=0.8pt, dash pattern=on 4pt off 3pt]
            ([xshift=0.55\textwidth+8pt]frame.north west) --
            ([xshift=0.55\textwidth+8pt]frame.south west);
    }
]
\textbf{Q.29} \quad {\footnotesize\textsc{[MSQ $\cdot$ 2 Mark(s) $\cdot$ No Neg]}} \hfill {\footnotesize\textbf{GATE 2025 CS2} \quad \tiny\texttt{2025\_CS2\_CS\_Q46}}
\vspace{2pt}\hrule height 0.3pt\vspace{4pt}

Consider the following relational schema along with all the functional dependencies that hold on them:
$R_1(A, B, C, D)$ with FDs, and $R_2(E, F, G, H)$ with FDs.
Which of the following statements is/are CORRECT?

\par\vspace{6pt}\noindent
\begin{tabular}{@{}p{0.48\linewidth}@{\hspace{0.04\linewidth}}p{0.48\linewidth}@{}}
$\square$\ \textbf{(A)} $R_1$ is in 3NF & $\square$\ \textbf{(B)} $R_2$ is in 3NF \\
\addlinespace[3pt]
$\square$\ \textbf{(C)} $R_1$ is NOT in 3NF & $\square$\ \textbf{(D)} $R_2$ is NOT in 3NF
\end{tabular}

\tcblower

\textbf{Official Key:} \textbf{\large C;D} \hfill \textit{\small [Medium]}
\vspace{2pt}\hrule height 0.3pt\vspace{3pt}
\textbf{Topic:} Databases \\
\textbf{Subtopic:} \textit{Integrity constraints, normal forms: 3NF and BCNF}
\vspace{2pt}\hrule height 0.3pt\vspace{3pt}
\textbf{Guide / Hint:} \\
{\small A schema is in 3NF if for every non-trivial FD $X \to Y$, either $X$ is a superkey or $Y$ consists of prime attributes. Here, both schemas contain dependencies where the LHS is not a superkey and the RHS has non-prime attributes, so neither is in 3NF.}

\end{tcolorbox}
\vspace{6pt}\hrule height 0.35pt\vspace{6pt}
\noindent
\begin{tcolorbox}[
    blanker,
    sidebyside,
    sidebyside align=top,
    lefthand width=0.55\textwidth,
    top=4pt, bottom=4pt, left=0pt, right=0pt,
    middle=16pt,
    underlay={
        \draw[dashed, black!80, line width=0.8pt, dash pattern=on 4pt off 3pt]
            ([xshift=0.55\textwidth+8pt]frame.north west) --
            ([xshift=0.55\textwidth+8pt]frame.south west);
    }
]
\textbf{Q.30} \quad {\footnotesize\textsc{[MSQ $\cdot$ 2 Mark(s) $\cdot$ No Neg]}} \hfill {\footnotesize\textbf{GATE 2025 CS2} \quad \tiny\texttt{2025\_CS2\_CS\_Q54}}
\vspace{2pt}\hrule height 0.3pt\vspace{4pt}

Consider the relational schema:
- $\text{Students}(\underline{\text{rollno}}, \text{name}, \text{age}, \text{cgpa})$
- $\text{Enrolled}(\underline{\text{rollno}, \text{courseno}})$
\par\medskip
Which of the following SQL query/queries correctly find(s) the name of all students who have enrolled in course 1470?

\par\vspace{6pt}\noindent
\begin{tabular}{@{}p{0.98\linewidth}@{}}
$\square$\ \textbf{(A)} SELECT S.name \newline FROM Students S \newline WHERE EXISTS (SELECT * FROM Enrolled E WHERE E.courseno = 1470 AND E.rollno = S.rollno); \\
\addlinespace[3pt]
$\square$\ \textbf{(B)} SELECT S.name \newline FROM Students S \newline WHERE SIZEOF (SELECT * FROM Enrolled E WHERE E.courseno = 1470 AND E.rollno = S.rollno) > 0; \\
\addlinespace[3pt]
$\square$\ \textbf{(C)} SELECT S.name \newline FROM Students S \newline WHERE 0 < (SELECT COUNT(*) FROM Enrolled E WHERE E.courseno = 1470 AND E.rollno = S.rollno); \\
\addlinespace[3pt]
$\square$\ \textbf{(D)} SELECT S.name \newline FROM Students S NATURAL JOIN Enrolled E \newline WHERE E.courseno = 1470;
\end{tabular}

\tcblower

\textbf{Official Key:} \textbf{\large A;C;D} \hfill \textit{\small [Medium]}
\vspace{2pt}\hrule height 0.3pt\vspace{3pt}
\textbf{Topic:} Databases \\
\textbf{Subtopic:} \textit{Relational model: SQL queries and joins}
\vspace{2pt}\hrule height 0.3pt\vspace{3pt}
\textbf{Guide / Hint:} \\
{\small Queries A, C, and D all correctly return the names of students enrolled in course 1470 using EXISTS, COUNT, and natural join respectively. Query B is invalid because SQL does not have a `SIZEOF` operator.}

\end{tcolorbox}
\vspace{6pt}\hrule height 0.35pt\vspace{6pt}
\noindent
\begin{tcolorbox}[
    blanker,
    sidebyside,
    sidebyside align=top,
    lefthand width=0.55\textwidth,
    top=4pt, bottom=4pt, left=0pt, right=0pt,
    middle=16pt,
    underlay={
        \draw[dashed, black!80, line width=0.8pt, dash pattern=on 4pt off 3pt]
            ([xshift=0.55\textwidth+8pt]frame.north west) --
            ([xshift=0.55\textwidth+8pt]frame.south west);
    }
]
\textbf{Q.31} \quad {\footnotesize\textsc{[NAT $\cdot$ 2 Mark(s) $\cdot$ No Neg]}} \hfill {\footnotesize\textbf{GATE 2025 CS2} \quad \tiny\texttt{2025\_CS2\_CS\_Q57}}
\vspace{2pt}\hrule height 0.3pt\vspace{4pt}

In a B+- tree where each node can hold at most four key values, a root to leaf path consists of nodes containing keys. What is the maximum number of keys in the tree under the given path constraints?

\par\vspace{6pt}\noindent\textbf{Your Answer:} \fbox{\rule{0pt}{1.3ex}\hspace{3cm}}

\tcblower

\textbf{Official Key:} \textbf{\large 33 to 33} \hfill \textit{\small [Hard]}
\vspace{2pt}\hrule height 0.3pt\vspace{3pt}
\textbf{Topic:} Databases \\
\textbf{Subtopic:} \textit{File organization, indexing: B and B+ trees}
\vspace{2pt}\hrule height 0.3pt\vspace{3pt}
\textbf{Guide / Hint:} \\
{\small In a B+ tree of order 5 (max 4 keys, max 5 children per internal node), min keys per non-root internal node is $\lceil 5/2 \rceil - 1 = 2$. Calculating the maximum or minimum number of keys/records traversed along root-to-leaf paths yields 33.}

\end{tcolorbox}
\vspace{6pt}\hrule height 0.35pt\vspace{6pt}
\end{document}